% GENERATED from content/data/enrichment.json by mxbuild.enrichment; do not edit.
\mxunit{enrichment-pigeonhole-notes}{Le principe des tiroirs}
\mxsubtitle{Seconde · olympiades \textperiodcentered{} Leçon et méthode \textperiodcentered{} 02}

\begin{prerequisites}
Division euclidienne et dénombrement élémentaire.
\end{prerequisites}

\begin{keypoint}[Idées et méthode]
Si $n + 1$ objets sont rangés dans n tiroirs, au moins un tiroir contient deux objets. Sinon, chaque tiroir en contiendrait au plus un et il y aurait au plus n objets.

Plus généralement, si N objets occupent k tiroirs, un tiroir en contient au moins $\lceil \frac{N}{k}\rceil $. Le plafond $\lceil x\rceil $ est le plus petit entier supérieur ou égal à x.

Le choix des tiroirs fait la preuve : restes d’une division, mois de naissance, petits carrés d’un quadrillage ou classes de nombres. Il faut justifier que chaque objet appartient à un tiroir.
\end{keypoint}

\begin{example}
Parmi 13 personnes, deux sont nées le même mois : les 12 mois sont les tiroirs. Cela ne garantit pas deux anniversaires le même jour.
\end{example}

\begin{activity}[À explorer]
Choisir six entiers distincts entre 1 et 10. Les grouper par paires de somme 11 et repérer une paire complète.
\end{activity}

